\section{Phân tích Use Case}
\label{sec:use-cases}

\subsection{Phạm vi và tác nhân}

BrewLite hỗ trợ khách hàng chọn đồ uống, tùy chỉnh món, quản lý giỏ hàng,
đăng nhập, tạo đơn và thanh toán không dùng tiền mặt để nhận tại quầy.
Phần này trình bày mô hình ca sử dụng của Assignment~1. Các đặc tả tập trung
vào hành vi quan sát được và điều kiện chấp nhận; việc có mã nguồn tương ứng
không đồng nghĩa mọi điều kiện đã được kiểm thử hoặc hệ thống đã sẵn sàng
vận hành thanh toán thật.

Luồng khách hàng UC-01--UC-07 là trọng tâm MVP. Quản trị đơn hàng ở UC-08,
quản lý nguyên liệu theo công thức và tích hợp ZaloPay Sandbox được ghi nhận
là phần mở rộng hiện có trong mã nguồn. Chúng giúp làm rõ cách thực thi
quy tắc nghiệp vụ, không thay thế phạm vi MVP bằng một hệ thống quản trị đầy đủ.

\begin{longtable}{|p{3.0cm}|p{11.6cm}|}
\caption{Tác nhân và trách nhiệm trong mô hình Use Case}
\label{tab:uc-actors}\\ \hline
\textbf{Tác nhân} & \textbf{Vai trò} \\ \hline
\endfirsthead
\hline \textbf{Tác nhân} & \textbf{Vai trò} \\ \hline
\endhead
Khách hàng & Xem menu khi chưa đăng nhập; chọn món và lưu giỏ trong trình
duyệt; xác thực trước khi tạo đơn, thanh toán hoặc xem dữ liệu đơn cá nhân. \\ \hline
Quản trị viên & Tác nhân của phần mở rộng: xem đơn của cửa hàng, cập nhật
tiến độ pha chế và hủy đơn còn đủ điều kiện. Phải được xác thực và có quyền
quản trị; không phải mọi khách hàng đã đăng nhập đều có quyền này. \\ \hline
Cổng thanh toán & Hệ thống bên ngoài cung cấp kết quả thanh toán. Trong
MVP, cổng giả lập đóng vai trò mô phỏng phụ thuộc này; phần mở rộng sử dụng
ZaloPay Sandbox. Không có giao dịch tiền thật trong luồng thử nghiệm. \\ \hline
\end{longtable}

Backend, cơ sở dữ liệu, bộ xác thực JWT và tác vụ định kỳ là thành phần
nội bộ của BrewLite, không được mô hình hóa thành tác nhân bên ngoài.

\subsection{Sơ đồ Use Case tổng quát}

\begin{figure}[htbp]
\centering
\begin{tikzpicture}[
  x=1cm,y=1cm,font=\small,
  uc/.style={ellipse,draw,align=center,text width=5.5cm,
    minimum height=0.96cm,inner sep=3pt},
  actorlabel/.style={align=center,text width=2.5cm,font=\small}
]
  \draw[rounded corners=2pt] (2.1,0.9) rectangle (10.8,-11.8);
  \node[font=\bfseries] at (6.45,0.5) {Hệ thống BrewLite};
  \node[uc] (uc01) at (6.45,-0.5) {UC-01: Xem và tìm món};
  \node[uc] (uc02) at (6.45,-2.0) {UC-02: Xem chi tiết, chọn tùy chọn};
  \node[uc] (uc03) at (6.45,-3.5) {UC-03: Quản lý giỏ hàng};
  \node[uc] (uc04) at (6.45,-5.0) {UC-04: Đăng ký, đăng nhập};
  \node[uc] (uc05) at (6.45,-6.5) {UC-05: Tạo đơn hàng};
  \node[uc] (uc06) at (6.45,-8.0) {UC-06: Thanh toán đơn hàng};
  \node[uc] (uc07) at (6.45,-9.5) {UC-07: Xem xác nhận và lịch sử};
  \node[uc,dashed] (uc08) at (6.45,-11.0) {UC-08: Xử lý trạng thái đơn\\(mở rộng quản trị)};
  % UML actors are drawn directly to keep the report self-contained.
  \draw (0.8,-4.1) circle (0.18);
  \draw (0.8,-4.28)--(0.8,-4.95);
  \draw (0.35,-4.55)--(1.25,-4.55);
  \draw (0.8,-4.95)--(0.4,-5.45) (0.8,-4.95)--(1.2,-5.45);
  \node[actorlabel] at (0.8,-5.8) {Khách hàng};
  \coordinate (customer) at (1.3,-4.55);
  \foreach \target in {uc01,uc02,uc03,uc04,uc05,uc06,uc07}
    \draw (customer)--(\target.west);
  \draw (12.0,-7.1) circle (0.18);
  \draw (12.0,-7.28)--(12.0,-7.95);
  \draw (11.55,-7.55)--(12.45,-7.55);
  \draw (12.0,-7.95)--(11.6,-8.45) (12.0,-7.95)--(12.4,-8.45);
  \node[actorlabel] at (12.0,-8.9) {Cổng thanh toán};
  \draw (11.55,-7.55)--(uc06.east);
  \draw (12.0,-9.8) circle (0.18);
  \draw (12.0,-9.98)--(12.0,-10.65);
  \draw (11.55,-10.25)--(12.45,-10.25);
  \draw (12.0,-10.65)--(11.6,-11.15) (12.0,-10.65)--(12.4,-11.15);
  \node[actorlabel] at (12.0,-11.65) {Quản trị viên};
  \draw (11.55,-10.25)--(uc08.east);
\end{tikzpicture}
\caption{Các mục tiêu sử dụng BrewLite và ranh giới hệ thống}
\label{fig:overall-use-case}
\end{figure}

Đường nối trong Hình~\ref{fig:overall-use-case} biểu diễn sự tham gia của
tác nhân, không biểu diễn thứ tự thao tác. Xác thực là tiền điều kiện của
các ca sử dụng được bảo vệ; khách hàng có phiên hợp lệ không cần thực hiện
lại UC-04 mỗi khi tạo đơn. Quy tắc trạng thái được dùng chung trong
UC-05--UC-08; UC-08 thể hiện riêng mục tiêu vận hành của quản trị viên.

\subsection{Danh mục và truy vết yêu cầu}

\begin{longtable}{|p{1.3cm}|p{3.6cm}|p{6.3cm}|p{2.7cm}|}
\caption{Danh mục Use Case và liên kết yêu cầu chức năng}
\label{tab:uc-catalog}\\ \hline
\textbf{Mã} & \textbf{Tên} & \textbf{Kết quả mong muốn} & \textbf{Yêu cầu} \\ \hline
\endfirsthead
\hline \textbf{Mã} & \textbf{Tên} & \textbf{Kết quả mong muốn} & \textbf{Yêu cầu} \\ \hline
\endhead
UC-01 & Xem và tìm món & Có danh sách món, giá, ảnh, danh mục và tình trạng
phục vụ; tìm kiếm, lọc theo danh mục. & FR-01 \\ \hline
UC-02 & Chi tiết và tùy chọn & Chọn size, topping được hỗ trợ, số lượng và
xem giá tạm tính trước khi thêm giỏ. & FR-02 \\ \hline
UC-03 & Quản lý giỏ hàng & Xem, sửa, xóa món; cập nhật tổng tạm tính và
chuyển sang đặt hàng. & FR-03 \\ \hline
UC-04 & Đăng ký, đăng nhập & Nhận phiên xác thực để thực hiện giao dịch và
truy cập dữ liệu cá nhân. & FR-04 \\ \hline
UC-05 & Tạo đơn hàng & Có mã đơn, giá đã xác nhận và trạng thái chờ thanh
toán; giữ tồn kho, áp dụng ưu đãi hợp lệ. & FR-05, FR-07, FR-09, FR-10 \\ \hline
UC-06 & Thanh toán đơn hàng & Ghi nhận kết quả thanh toán nhất quán,
không tạo giao dịch trùng khi gửi lại cùng yêu cầu. & FR-05, FR-07, FR-08, FR-10 \\ \hline
UC-07 & Xác nhận và lịch sử & Xem mã nhận món, nội dung đơn, tổng tiền,
trạng thái và lịch sử đơn của bản thân. & FR-06, FR-07 \\ \hline
UC-08 & Xử lý trạng thái đơn & Quản trị viên cập nhật tiến độ hợp lệ và có
dấu vết thay đổi; hủy đơn theo quy tắc. & FR-07, FR-09, FR-10; mở rộng \\ \hline
\end{longtable}

Trong bảng trên, FR-07 là máy trạng thái đơn; FR-08 là tính lũy đẳng
thanh toán; FR-09 là xử lý tồn kho khi có truy cập đồng thời; FR-10 là
khuyến mãi và tích điểm. Các yêu cầu này là quy tắc xuyên suốt nhiều ca
sử dụng, không buộc khách hàng phải thao tác với một màn hình riêng.

\subsection{Đặc tả UC-01: Xem và tìm món}
\label{uc:01}
\label{uc01}

\begin{longtable}{|p{3.1cm}|p{11.5cm}|}
\hline \textbf{Thuộc tính} & \textbf{Nội dung} \\ \hline
Tác nhân, kích hoạt & Khách hàng mở trang menu hoặc thay đổi từ khóa,
danh mục. \\ \hline
Tiền điều kiện & Truy cập được ứng dụng; không bắt buộc đăng nhập. \\ \hline
Hậu điều kiện & Hiển thị danh sách phù hợp hoặc trạng thái rỗng/lỗi rõ
ràng; chưa tạo đơn hay thay đổi tồn kho. \\ \hline
Liên kết & FR-01; chuyển sang UC-02 khi chọn món. \\ \hline
\end{longtable}

\textbf{Luồng chính}
\begin{enumerate}[nosep]
  \item Khách hàng mở \path{/menu}; hệ thống hiển thị trạng thái đang tải.
  \item Hệ thống lấy danh sách qua \path{GET /api/v1/products} và hiển thị
  tên, ảnh, giá cơ bản, danh mục, tình trạng phục vụ.
  \item Khách hàng nhập từ khóa hoặc chọn danh mục.
  \item Giao diện lọc kết hợp hai điều kiện trên danh sách đã tải; khách
  hàng chọn món để xem chi tiết hoặc mở hộp chọn tùy chọn.
\end{enumerate}

\textbf{Thay thế và ngoại lệ.} Danh sách chưa có món được hiển thị bằng
trạng thái rỗng. Không có kết quả tìm kiếm/lọc được thông báo riêng. Khi
lỗi mạng, giao diện báo lỗi và cho phép thử lại. Món không khả dụng vẫn
có thể xuất hiện với nhãn hết hàng và nút mua bị khóa; không suy diễn rằng
mọi món hiển thị đều chắc chắn còn nguyên liệu tại thời điểm đặt đơn.

\textbf{Điều kiện chấp nhận.} Kết quả kết hợp đúng từ khóa và danh mục;
tìm tên không phân biệt hoa/thường và bỏ khoảng trắng đầu/cuối. Mã hiện
tại chưa chuẩn hóa dấu tiếng Việt, vì vậy tìm không dấu chưa được coi là
tiêu chí đã đáp ứng. Nút mua không hoạt động với món hết hàng hoặc chưa có
size khả dụng.

\subsection{Đặc tả UC-02: Xem chi tiết và chọn tùy chọn}
\label{uc:02}
\label{uc02}

\begin{longtable}{|p{3.1cm}|p{11.5cm}|}
\hline \textbf{Thuộc tính} & \textbf{Nội dung} \\ \hline
Tác nhân, kích hoạt & Khách hàng chọn một món trong menu hoặc mở trang
chi tiết món. \\ \hline
Tiền điều kiện & Có định danh sản phẩm hợp lệ; để thêm giỏ, món và size
phải đang khả dụng. \\ \hline
Hậu điều kiện & Dòng món với cấu hình đã chọn được lưu vào giỏ của trình
duyệt; nếu thao tác thất bại, không thêm dòng không hợp lệ. \\ \hline
Liên kết & FR-02; dữ liệu được dùng tiếp bởi UC-03 và xác minh lại tại UC-05. \\ \hline
\end{longtable}

\textbf{Luồng chính}
\begin{enumerate}[nosep]
  \item Hệ thống lấy chi tiết qua \path{GET /api/v1/products/:id}.
  \item Khách hàng xem thông tin món và chọn size trong các size được hỗ trợ.
  \item Với đồ uống ngoài nhóm cà phê, khách hàng có thể chọn tối đa ba
  loại topping khả dụng, không trùng lặp. Cà phê chỉ chọn size, không dùng topping.
  \item Khách hàng chọn số lượng từ 1 đến 99; giao diện cập nhật giá tạm tính.
  \item Khách hàng xác nhận; hệ thống thêm cấu hình vào giỏ và thông báo kết quả.
\end{enumerate}

\textbf{Thay thế và ngoại lệ.} Khi sửa một dòng trong giỏ, cấu hình mới
thay thế dòng được chọn. Không có topping khả dụng thì khách hàng vẫn có
thể mua món không topping. Món không tồn tại, lỗi tải, size hết hàng hoặc
dữ liệu giá không hợp lệ phải được thông báo; không cho thêm món khi chưa
có cấu hình hợp lệ. Nếu nguyên liệu thay đổi sau lúc xem chi tiết, kết
quả cuối cùng được quyết định tại UC-05.

\textbf{Điều kiện chấp nhận.} Giá hiển thị cho một dòng bằng số lượng
nhân với tổng giá cơ bản, phụ thu size và giá topping. Cấu hình cà phê
không chứa topping; các giới hạn số lượng và topping được kiểm tra lại ở
backend khi tạo đơn. Giá trong trình duyệt là ước tính, không phải giá
được backend tin cậy.

\subsection{Đặc tả UC-03: Quản lý giỏ hàng}
\label{uc:03}
\label{uc03}

\begin{longtable}{|p{3.1cm}|p{11.5cm}|}
\hline \textbf{Thuộc tính} & \textbf{Nội dung} \\ \hline
Tác nhân, kích hoạt & Khách hàng mở giỏ, sửa cấu hình, tăng/giảm số lượng
hoặc xóa một món. \\ \hline
Tiền điều kiện & Có thể mở giỏ khi chưa đăng nhập; có dòng món nếu muốn
sửa, xóa hoặc tiếp tục đặt hàng. \\ \hline
Hậu điều kiện & Giỏ trong trình duyệt và tổng tạm tính được cập nhật;
chưa trừ tồn kho và chưa ghi nhận thanh toán. \\ \hline
Liên kết & FR-03; UC-04 nếu cần đăng nhập, sau đó UC-05. \\ \hline
\end{longtable}

\textbf{Luồng chính}
\begin{enumerate}[nosep]
  \item Khách hàng mở \path{/cart}; hệ thống nạp giỏ đã lưu trong trình duyệt.
  \item Hệ thống trình bày tên món, size, topping, số lượng và thành tiền tạm tính.
  \item Khách hàng chỉnh số lượng hoặc cấu hình; các dòng cùng sản phẩm,
  size và cùng tập topping được gộp theo quy tắc của giỏ.
  \item Hệ thống lưu thay đổi và tính lại tổng tiền.
  \item Khách hàng xác nhận các món để chuyển sang bước đặt hàng.
\end{enumerate}

\textbf{Thay thế và ngoại lệ.} Khách hàng xóa dòng không còn cần; xóa hết
thì hiển thị giỏ rỗng và không tiếp tục checkout. Số lượng ngoài giới hạn
không được chấp nhận. Nếu chưa đăng nhập, khách hàng được yêu cầu xác
thực trước khi đặt hàng. Giỏ hiện tại của giao diện khách hàng dùng
\texttt{localStorage}; các API Cart có trong backend không phải bằng
chứng về đồng bộ giỏ đa thiết bị của giao diện này.

\textbf{Điều kiện chấp nhận.} Sửa/xóa cập nhật ngay giỏ và tổng tiền; hai
cấu hình khác size hoặc khác tập topping không bị gộp nhầm. Không cho gửi
đơn với giỏ rỗng; backend giới hạn tối đa 50 dòng trong một yêu cầu tạo đơn.

\subsection{Đặc tả UC-04: Đăng ký và đăng nhập}
\label{uc:04}
\label{uc04}

\begin{longtable}{|p{3.1cm}|p{11.5cm}|}
\hline \textbf{Thuộc tính} & \textbf{Nội dung} \\ \hline
Tác nhân, kích hoạt & Khách hàng chọn đăng ký/đăng nhập hoặc cần thực hiện
một thao tác được bảo vệ. \\ \hline
Tiền điều kiện & Chưa có phiên hợp lệ; để đăng nhập cần tài khoản đã tồn tại. \\ \hline
Hậu điều kiện & Thành công: có JWT và thông tin người dùng. Thất bại:
không cấp phiên cho thông tin sai hoặc dữ liệu không hợp lệ. \\ \hline
Liên kết & FR-04; tiền điều kiện của UC-05--UC-07 và truy cập quản trị. \\ \hline
\end{longtable}

\textbf{Luồng chính: đăng nhập}
\begin{enumerate}[nosep]
  \item Khách hàng mở biểu mẫu, nhập email và mật khẩu.
  \item Giao diện gửi \path{POST /api/v1/auth/login}.
  \item Backend kiểm tra dữ liệu, đối chiếu mật khẩu với giá trị đã băm.
  \item Hệ thống cấp JWT, trả thông tin tài khoản; giao diện cập nhật phiên.
  \item Khách hàng tiếp tục thao tác cần xác thực.
\end{enumerate}

\textbf{Luồng thay thế: đăng ký.} Khách hàng nhập họ tên, email, mật khẩu
và xác nhận mật khẩu. Giao diện gửi \path{POST /api/v1/auth/register};
backend kiểm tra dữ liệu và email duy nhất, lưu mật khẩu đã băm rồi cấp
phiên cho tài khoản mới. Điểm tích lũy ban đầu bằng không.

\textbf{Ngoại lệ.} Email đã dùng được báo xung đột; dữ liệu sai định dạng
được báo để sửa. Khi đăng nhập sai, hệ thống trả thông báo chung về thông
tin đăng nhập. JWT thiếu, hết hạn hoặc không hợp lệ không được truy cập
API bảo vệ; giao diện cần yêu cầu đăng nhập lại.

\textbf{Điều kiện chấp nhận.} Không trả mật khẩu hay giá trị băm cho client.
Đơn hàng cá nhân luôn gắn với người dùng được xác thực, không nhận danh
tính chủ đơn tùy ý từ trình duyệt. Có phiên đăng nhập không tự động cấp
quyền quản trị.

\subsection{Đặc tả UC-05: Tạo đơn hàng}
\label{uc:05}
\label{uc05}

\begin{longtable}{|p{3.1cm}|p{11.5cm}|}
\hline \textbf{Thuộc tính} & \textbf{Nội dung} \\ \hline
Tác nhân, kích hoạt & Khách hàng xác nhận đặt hàng tại trang checkout. \\ \hline
Tiền điều kiện & Phiên đăng nhập hợp lệ; giỏ từ 1 đến 50 dòng; mỗi dòng
chứa sản phẩm, size, topping và số lượng hợp lệ. \\ \hline
Hậu điều kiện & Thành công: có đơn \texttt{PENDING}, mã đơn, chi tiết và
giá đã xác nhận; nguyên liệu/topping được giữ bằng cách trừ kho, điểm sử
dụng được ghi nhận. Thất bại: giao dịch rollback, không lưu đơn dở dang
hoặc trừ kho/điểm một phần. \\ \hline
Liên kết & FR-05, FR-07, FR-09, FR-10; tiếp tục UC-06. \\ \hline
\end{longtable}

\textbf{Luồng chính}
\begin{enumerate}[nosep]
  \item Khách hàng xem món đã chọn, có thể nhập mã khuyến mãi và số điểm muốn dùng.
  \item Giao diện kiểm tra đơn đang chờ gắn với giỏ hiện tại để có thể tiếp
  tục đơn trước đó, sau đó gửi \path{POST /api/v1/orders} khi cần tạo đơn mới.
  Yêu cầu chỉ chứa cấu hình món và lựa chọn ưu đãi, không chứa giá được tin cậy.
  \item Backend xác thực người dùng, kiểm tra sản phẩm đang bán, size hỗ
  trợ, topping cho phép và tính lại toàn bộ giá từ cơ sở dữ liệu.
  \item Backend cộng nhu cầu nguyên liệu từ công thức theo sản phẩm, size
  và số lượng; kiểm tra và trừ nguyên liệu cùng topping trong giao dịch.
  \item Backend kiểm tra hiệu lực, điều kiện mã khuyến mãi và số dư điểm;
  tính số tiền cuối cùng, tạo đơn \texttt{PENDING}, lưu chi tiết và lịch sử trạng thái.
  \item Hệ thống trả mã đơn và số tiền đã xác nhận; giao diện lưu tham
  chiếu đơn đang chờ rồi chuyển đến trang thanh toán.
\end{enumerate}

\textbf{Luồng thay thế}
\begin{itemize}[nosep]
  \item Giỏ chưa thay đổi và đã có đơn \texttt{PENDING} hoặc
  \texttt{PAYMENT\_FAILED}: giao diện tiếp tục đơn đó. Nếu giỏ đã sửa, giao
  diện xử lý hủy đơn chờ cũ trước khi tạo đơn mới theo cơ chế hiện có.
  \item Không nhập mã hay không dùng điểm: tính đơn theo giá thông thường.
  Mã hợp lệ được áp dụng trước phần giảm bằng điểm.
  \item Điểm sử dụng là bội số của 10, mỗi điểm giảm 100~đồng; không vượt
  số dư và số tiền sau khuyến mãi. Điều kiện này được xác minh tại backend.
\end{itemize}

\textbf{Ngoại lệ.} Từ chối giỏ rỗng, sản phẩm không bán, size không có
công thức, topping không hợp lệ, mã ưu đãi không đủ điều kiện hoặc điểm
không đủ. Khi thiếu nguyên liệu hoặc dữ liệu tồn kho thay đổi đồng thời,
trả thông báo để khách tải lại/thử lại; không âm kho và không ghi đơn một phần.

\textbf{Điều kiện chấp nhận.} Tổng tiền do backend quyết định. Tồn kho
đồ uống được tính từ nguyên liệu của công thức theo size, không mô tả
bằng một trường tồn kho sản phẩm độc lập. Cập nhật nguyên liệu kiểm tra
\texttt{version} và số dư, kết hợp giao dịch \texttt{Serializable}; topping
được trừ có điều kiện đủ kho. Khi hai khách tranh lượng nguyên liệu cuối,
chỉ yêu cầu đáp ứng đủ điều kiện mới thành công, yêu cầu còn lại nhận lỗi
có thể xử lý. Cơ chế tiếp tục đơn ở giao diện không được suy rộng thành
cam kết lũy đẳng cho mọi yêu cầu tạo đơn từ client bất kỳ.

\subsection{Đặc tả UC-06: Thanh toán đơn hàng}
\label{uc:06}
\label{uc06}

\begin{longtable}{|p{3.1cm}|p{11.5cm}|}
\hline \textbf{Thuộc tính} & \textbf{Nội dung} \\ \hline
Tác nhân, kích hoạt & Khách hàng chọn phương thức và xác nhận thanh toán;
cổng thanh toán cung cấp kết quả. \\ \hline
Tiền điều kiện & JWT hợp lệ; đơn thuộc khách hàng, đang ở
\texttt{PENDING} hoặc \texttt{PAYMENT\_FAILED}; không có giao dịch thành
công hoặc giao dịch đang xử lý khác chặn lần thanh toán này. \\ \hline
Hậu điều kiện & Thành công: Payment là \texttt{SUCCESS}, đơn là
\texttt{PAID}, điểm thưởng và lịch sử được ghi nhất quán. Thất bại:
Payment là \texttt{FAILED}, đơn là \texttt{PAYMENT\_FAILED}; có thể thử
lại bằng yêu cầu mới theo quy tắc. \\ \hline
Liên kết & FR-05, FR-07, FR-08, FR-10; kết quả được xem tại UC-07. \\ \hline
\end{longtable}

\textbf{Luồng chính của thanh toán giả lập trong MVP}
\begin{enumerate}[nosep]
  \item Hệ thống nạp đơn từ backend, hiển thị số tiền đã xác nhận; khách
  hàng chọn ví điện tử hoặc thẻ nội địa trong luồng mô phỏng.
  \item Giao diện gửi \path{POST /api/v1/payments} với mã đơn, phương thức
  và header \texttt{Idempotency-Key} cho lần thanh toán.
  \item Backend kiểm tra quyền sở hữu, trạng thái đơn, khóa lũy đẳng và
  dấu vân tay nội dung yêu cầu; số tiền lấy từ đơn trong cơ sở dữ liệu.
  \item Backend tuần tự hóa các lần thanh toán của cùng đơn bằng khóa
  hàng, tạo giao dịch và nhận kết quả từ cổng giả lập.
  \item Với kết quả thành công, backend cập nhật giao dịch và chuyển đơn
  sang \texttt{PAID} trong cùng giao dịch dữ liệu; ghi lịch sử và cộng
  $\lfloor\text{số tiền cuối cùng}/1000\rfloor$ điểm thưởng.
  \item Giao diện hiển thị kết quả. Chỉ xóa giỏ tương ứng với đơn đang
  chờ khi giỏ hiện tại còn khớp bản đã đặt, tránh xóa các món vừa sửa.
\end{enumerate}

\textbf{Luồng thay thế và ngoại lệ}
\begin{itemize}[nosep]
  \item Cổng giả lập trả thất bại: lưu \texttt{FAILED}, chuyển đơn sang
  \texttt{PAYMENT\_FAILED}; giữ thông tin để khách thử lại hoặc hủy đơn.
  Lần thử mới dùng khóa mới; khi thử lại đơn chuyển về \texttt{PENDING}.
  \item Gửi lại cùng khóa và cùng nội dung sau khi đã hoàn tất: trả kết
  quả đã lưu, không tạo Payment và không cộng điểm lần nữa. Khóa trùng
  nhưng nội dung khác bị từ chối với xung đột.
  \item Mất phản hồi hoặc còn đang xử lý: client giữ khóa của lần gửi
  đó để kiểm tra/gửi lại, không suy ra thất bại chỉ vì chưa thấy kết quả.
  Giao dịch cùng khóa còn \texttt{PENDING} có thể trả xung đột yêu cầu thử lại.
  \item Sai chủ đơn, đơn đã hủy/đã thanh toán hoặc có giao dịch khác đang
  xử lý: từ chối, không tạo lần thu tiền mới.
\end{itemize}

\textbf{Điều kiện chấp nhận.} Gửi lặp hoặc gửi đồng thời không ghi nhận
hai lần thanh toán thành công cho một đơn. Giao dịch và trạng thái đơn
không mâu thuẫn; điểm chỉ được thưởng khi chuyển hợp lệ sang
\texttt{PAID}. ZaloPay Sandbox chỉ hoạt động khi \texttt{NODE\_ENV} là
\texttt{development}/\texttt{test}, \texttt{ZALOPAY\_SANDBOX\_ENABLED=true}
và có đủ bộ khoá sandbox; không dùng tiền thật.

\textbf{Luồng ZaloPay Sandbox.} Mã nguồn có luồng tạo yêu cầu
\path{POST /api/v1/payments/zalopay}, nhận callback và đồng bộ kết quả.
Luồng này có trạng thái chờ bất đồng bộ; callback phải được kiểm tra
chữ ký và dữ liệu giao dịch trước khi cập nhật đơn. Giao diện kiểm tra
giao dịch Sandbox đang chờ trước khi cho tạo lượt mới. Cấu hình sandbox,
khả năng nhận callback và kiểm thử tích hợp là điều kiện triển khai riêng;
báo cáo không coi đây là thanh toán thật đã nghiệm thu.

\subsection{Đặc tả UC-07: Xem xác nhận và lịch sử đơn}
\label{uc:07}
\label{uc07}

\begin{longtable}{|p{3.1cm}|p{11.5cm}|}
\hline \textbf{Thuộc tính} & \textbf{Nội dung} \\ \hline
Tác nhân, kích hoạt & Khách hàng xem kết quả sau thanh toán, mở danh sách
đơn hoặc chọn một đơn cụ thể. \\ \hline
Tiền điều kiện & Phiên hợp lệ; đơn cần xem thuộc khách hàng hiện tại. \\ \hline
Hậu điều kiện & Hiển thị dữ liệu đơn được backend xác nhận; việc xem
không tự đổi trạng thái hoặc tạo giao dịch mới. \\ \hline
Liên kết & FR-06, FR-07; có thể quay lại UC-06 nếu đơn còn được thanh toán. \\ \hline
\end{longtable}

\textbf{Luồng chính}
\begin{enumerate}[nosep]
  \item Khách hàng mở \path{/orders}; hệ thống gọi \path{GET /api/v1/orders/me}.
  \item Hệ thống hiển thị danh sách đơn cá nhân, mã đơn, thời gian, số
  lượng món, tổng tiền và trạng thái; hỗ trợ phân trang và lọc trạng thái.
  \item Khách hàng chọn đơn; hệ thống gọi \path{GET /api/v1/orders/:id}.
  \item Giao diện trình bày món và tùy chọn, các khoản tiền, kết quả
  thanh toán, lịch sử trạng thái và mã đơn dùng khi nhận tại quầy.
\end{enumerate}

\textbf{Thay thế và ngoại lệ.} Chưa có đơn thì hiển thị lịch sử rỗng.
Đơn chưa thanh toán có thể tiếp tục UC-06. Đơn đủ điều kiện có thể được
khách hàng hủy bằng \path{POST /api/v1/orders/:id/cancel}; quy tắc hoàn
kho, điểm và hoàn tiền theo UC-08. Phiên hết hạn yêu cầu đăng nhập lại;
đơn không tồn tại hoặc không thuộc khách được báo không tìm thấy, không
trả dữ liệu đơn người khác. Lỗi tải có thông báo và khả năng thử lại.

\textbf{Điều kiện chấp nhận.} Danh sách và chi tiết chỉ chứa đơn thuộc
người được xác thực. Nhãn thành công dựa trên trạng thái do backend trả
về, không dựa riêng vào trang chuyển hướng của cổng thanh toán. Không
khẳng định có hóa đơn email, mã QR nhận món hay thông báo đẩy nếu chưa
có luồng tương ứng trong phạm vi.

\subsection{Đặc tả UC-08: Xử lý trạng thái đơn hàng (mở rộng)}
\label{uc:08}
\label{uc08}

\begin{longtable}{|p{3.1cm}|p{11.5cm}|}
\hline \textbf{Thuộc tính} & \textbf{Nội dung} \\ \hline
Tác nhân, kích hoạt & Quản trị viên mở danh sách đơn và xác nhận tiến độ
thực tế tại quầy. \\ \hline
Tiền điều kiện & Phiên hợp lệ, có quyền quản trị; đơn tồn tại và chuyển
trạng thái được phép. \\ \hline
Hậu điều kiện & Trạng thái, phiên bản và lịch sử đơn được cập nhật nhất
quán; thay đổi sai quy tắc bị từ chối. Khi hủy, hoàn kho/điểm theo quy tắc. \\ \hline
Liên kết & FR-07, FR-09, FR-10; phần quản trị là mở rộng của MVP. \\ \hline
\end{longtable}

\textbf{Luồng chính}
\begin{enumerate}[nosep]
  \item Quản trị viên mở danh sách/chi tiết đơn qua các API
  \path{/api/v1/admin/orders} được bảo vệ bằng JWT và quyền quản trị.
  \item Quản trị viên chọn hành động tương ứng với công việc đã thực hiện.
  \item Hệ thống kiểm tra trạng thái và phiên bản hiện tại, chỉ chấp nhận
  các bước \texttt{PAID} $\rightarrow$ \texttt{PREPARING}
  $\rightarrow$ \texttt{READY} $\rightarrow$ \texttt{COMPLETED}.
  \item Hệ thống lưu trạng thái mới cùng người thay đổi, thời điểm và lý
  do; giao diện khách hàng có thể đọc trạng thái mới khi tải lại dữ liệu.
\end{enumerate}

\textbf{Luồng thay thế: hủy đơn.} Khách hàng sở hữu đơn hoặc quản trị
viên được hủy khi đơn ở \texttt{PENDING}, \texttt{PAYMENT\_FAILED} hoặc
\texttt{PAID}, trước khi bắt đầu pha chế. Hệ thống chuyển sang
\texttt{CANCELLED}, hoàn đúng lượng nguyên liệu đã ghi với đơn, hoàn
topping và điểm đã sử dụng. Với đơn đã trả tiền, thu hồi điểm đã thưởng
và ghi \texttt{MANUAL\_REQUIRED}: tiền cần được hoàn thủ công, không có
cam kết hoàn tiền tự động. Gọi hủy lại đơn đã hủy không được hoàn kho/điểm lần hai.

\textbf{Tác vụ nội bộ hiện có.} Scheduler mỗi phút tìm đơn
\texttt{PENDING} được tạo hơn 15 phút để gọi cơ chế hủy và hoàn kho.
Quy tắc hiện tại không tự áp dụng cho \texttt{PAYMENT\_FAILED}; không
diễn giải thành việc hủy mọi đơn sau 15 phút, hay giữ món 15 phút sau thanh toán.

\textbf{Ngoại lệ và điều kiện chấp nhận.} Từ chối tài khoản không có
quyền, bỏ bước pha chế, quay lui trạng thái kết thúc hoặc hủy khi đã
\texttt{PREPARING}. Khi có cập nhật đồng thời, điều kiện phiên bản phải
ngăn ghi đè; người dùng nhận thông báo tải lại. Hủy phải qua nghiệp vụ
hủy có hoàn kho/điểm, không chỉ gán trực tiếp trạng thái.

\subsection{Quy tắc trạng thái và kiểm chứng chấp nhận}

\begin{longtable}{|p{3.1cm}|p{5.8cm}|p{5.2cm}|}
\caption{Các chuyển trạng thái được phép}
\label{tab:order-transitions}\\ \hline
\textbf{Trạng thái nguồn} & \textbf{Trạng thái đích} & \textbf{Sự kiện} \\ \hline
\endfirsthead
\hline \textbf{Trạng thái nguồn} & \textbf{Trạng thái đích} & \textbf{Sự kiện} \\ \hline
\endhead
Chưa có đơn & \texttt{PENDING} & Tạo đơn hợp lệ. \\ \hline
\texttt{PENDING} & \texttt{PAID}, \texttt{PAYMENT\_FAILED},
\texttt{CANCELLED} & Thanh toán thành công, thất bại hoặc hủy hợp lệ. \\ \hline
\texttt{PAYMENT\_FAILED} & \texttt{PENDING}, \texttt{CANCELLED} & Thử
thanh toán lại hoặc hủy. \\ \hline
\texttt{PAID} & \texttt{PREPARING}, \texttt{CANCELLED} & Bắt đầu pha
chế hoặc hủy trước pha chế; hủy cần hoàn tiền thủ công. \\ \hline
\texttt{PREPARING} & \texttt{READY} & Pha chế xong. \\ \hline
\texttt{READY} & \texttt{COMPLETED} & Đã giao đồ uống tại quầy. \\ \hline
\texttt{COMPLETED}, \texttt{CANCELLED} & Không có & Trạng thái kết thúc. \\ \hline
\end{longtable}

Những tình huống sau là cơ sở thiết kế kiểm thử và nghiệm thu, không phải
báo cáo kết quả chạy kiểm thử:
\begin{enumerate}[nosep]
  \item \textbf{Luồng mua hoàn chỉnh:} xem menu, chọn size/topping hợp
  lệ, thêm giỏ, đăng nhập, tạo đơn, thanh toán thành công và xem lại đúng
  đơn trong lịch sử (UC-01--UC-07; FR-01--FR-06).
  \item \textbf{Chuyển trạng thái sai:} thử chuyển
  \texttt{PAID} thẳng sang \texttt{READY}; hệ thống từ chối và giữ lịch
  sử nhất quán (UC-08; FR-07).
  \item \textbf{Thanh toán lặp:} gửi đồng thời cùng khóa/cùng nội dung;
  chỉ có một kết quả nghiệp vụ. Thử lại khóa đó với nội dung khác phải
  bị từ chối (UC-06; FR-08).
  \item \textbf{Tranh tồn kho:} hai khách cùng đặt lượng nguyên liệu
  cuối; kiểm tra không âm kho, đơn lỗi không để lại dữ liệu một phần và
  thông báo cho phép xử lý tiếp (UC-05; FR-09).
  \item \textbf{Ưu đãi và điểm:} mã hết hạn, chưa đạt mức tối thiểu,
  điểm vượt số dư hoặc không là bội số của 10 đều bị từ chối. Thanh toán
  lặp không thưởng điểm hai lần; hủy lặp không hoàn điểm hai lần
  (UC-05, UC-06, UC-08; FR-10).
  \item \textbf{Quyền truy cập:} khách A không xem/thanh toán đơn của
  khách B; tài khoản khách không cập nhật trạng thái qua API quản trị
  (UC-04, UC-06--UC-08).
\end{enumerate}

\subsection{Căn cứ và giới hạn của đặc tả}

Các ca sử dụng được tổng hợp từ đề Assignment~1
\path{BrewLite-CongNghePhanMem.pdf}, tài liệu các phân hệ trong
\path{docs/product-catalog}, \path{docs/product-detail}, \path{docs/cart},
\path{docs/auth}, \path{docs/order}, \path{docs/payment} và
\path{docs/order-history}. Hành vi cụ thể được đối chiếu với
\path{backend/src/modules/orders/orders.service.ts},
\path{backend/src/modules/orders/order-state.ts},
\path{backend/src/modules/products/products.service.ts},
\path{backend/src/modules/payments/zalopay-sandbox.service.ts},
\path{backend/src/modules/payments/zalopay-sandbox.gateway.ts},
\path{frontend/src/lib/pending-order.ts} và các trang giao diện liên quan.

Một số tài liệu phân hệ phản ánh thiết kế ở thời điểm cũ, ví dụ chỉ hiển
thị món còn hàng, tìm không dấu hoặc chưa có tác vụ hết hạn đơn. Khi mô
tả hiện trạng, phần này dùng hành vi trong mã nguồn đã đối chiếu và nêu
rõ phần mở rộng. Việc nghiệm thu cần thực thi các tình huống chấp nhận
trong môi trường cấu hình tương ứng, đặc biệt với xử lý đồng thời,
callback Sandbox và hoàn tiền thủ công.
